% 日本語論文。LuaTeX が必須で、pdflatex では日本語が PDF から丸ごと落ちる。
% 同梱テンプレート（iclr2024 / mdpi / agents4science_2025）は英語会議用で
% CJK に対応しないため使わず、プリアンブルを自前で書いている。
\documentclass[11pt]{article}
\usepackage{luatexja-fontspec}
\setmainjfont{IPAexMincho}
\setsansjfont{IPAexGothic}

\usepackage{amsmath}
\usepackage{amssymb}
\usepackage{graphicx}
\usepackage{booktabs}
\usepackage{caption}
\usepackage[margin=25mm]{geometry}
\usepackage[hidelinks]{hyperref}
\usepackage{natbib}
\bibliographystyle{unsrtnat}

\title{固定推論予算下における協調フォールディングの配分戦略：\\
MSA 多様化は届かず、別系統への配分が最も近づいた}
\author{AIRAS 自動研究パイプライン}
\date{2026年8月}

\begin{document}
\maketitle

\begin{abstract}
協調フォールディング（co-folding）手法は各系につき25構造を生成するが、その25枠を
どう配分すべきかは検証されていない。先行研究は多重配列アラインメント（MSA）の多様化が
精度を上げると報告するものの、いずれも多様化と同時に生成数を増やしており、多様性の
効果と計算量の効果を分離できていない。本研究は総数を25に固定した対照を置くことで
この交絡を除き、Runs N' Poses ベンチマークの205系で3つの配分を対比較した。
MSA 深度のサブサンプリングは区別できるポーズモード数を3.76から5.21へ増やし、摂動自体は
確実に作動したが、主指標 oracle@25 lDDT-PLI の対応差は $+0.0121$
（95\%区間 $[+0.0014, +0.0226]$）にとどまり、\textbf{事前に登録した閾値 $+0.030$ に
届かなかった}。効果は学習データ類似度に依存し、最も難しい 20--30 帯では負に転じた。
一方、予算の一部を別のアーキテクチャ系統（拡散ドッキング）に振り向ける配分は
$+0.0298$ $[+0.0187, +0.0425]$ を示し、\textbf{4つの類似度帯すべてで区間が0を跨がなかった}
ものの、平均は閾値を $0.0002$ 下回った。\textbf{我々はこの系統混合について事前に負の結果を
期待値として登録しており、予測は外れた}。ドッキングは予算の20\%を占め、最良ポーズを
供給した割合も20.0\%と比例しており、1本あたりの質ではなく分布の非相関性が効いている。
なお3アームすべてで、候補集合が広いほど信頼度による選択は悪化した。本研究が主張するのは
固定推論予算下での配分戦略の比較であって、新しい協調フォールディングモデルではない。
\end{abstract}

\section{はじめに}

タンパク質とリガンドの複合体構造を配列と SMILES だけから予測する協調フォールディング
手法は、AlphaFold3 を含む複数の実装が公開され、実用段階にあるとされてきた。
しかし \citet{skrinjar-2026-have} は、学習カットオフ後の2,600系からなるベンチマーク
Runs N' Poses を構築し、現行手法が学習データ中のリガンド配座を大きく記憶していること、
そして「サンプリングとランキングの両方に改善が必要」であることを示した。とりわけ
著者らは、\textbf{正解ポーズをサンプリングすること自体が主ボトルネック}であり、
かつ\textbf{推論時に複数シードを使っても有用でない}と報告している。

ここに、これまで測られていない問いがある。25構造という推論予算そのものは所与として、
\textbf{その配り方を変えたらどうなるのか}。シードに配っても多様性が生まれないなら、
より上流を揺らせばよいはずである。

MSA を摂動する操作が AlphaFold 系列のサンプリングを改善することは複数報告されている。
\citet{lidbrink-2026-several} は確率的サブサンプリング・クラスタリング・カラムマスクの
3手法が AlphaFold3 の代替立体構造状態のサンプリングを統計的に有意に改善することを
107タンパク質で示した。\citet{neupane-2026-improving} は MSA とテンプレートの多様化が
タンパク質-リガンド複合体の ligand RMSD を改善すると報告している。

\textbf{しかしいずれも、多様化と同時にモデル生成数を増やしている。}前者は各手法
1,000構造以上を生成し、後者は MSA--テンプレート組合せごとに25--100モデルを生成して
プールする。この設計では「多様性が効いた」のか「単により多く計算したから当たった」
のかを原理的に区別できない。実務家が直面するのは常に固定予算の配分問題であり、
この区別こそが意思決定にとって本質的である。

本研究は総構造数を25に固定した対照を置くことでこの交絡を設計段階で消す。

\section{関連研究}

Runs N' Poses \citep{skrinjar-2026-have} は本研究のベンチマークであり、
その取り込みパイプラインは PLINDER \citep{durairaj-2024-plinder} に由来する。
本研究の採点設定も plinder.eval に準拠する。

MSA 摂動については前述の2件 \citep{lidbrink-2026-several,neupane-2026-improving} に加え、
\citet{xing-2025-state} が状態特異的な MSA 部分集合の選択によりリガンドポーズ予測を
改善する事例を報告している。\textbf{機構としての「MSA を振る」は既に確立されており、
本研究の新規性はそこにはない。}本研究が加えるのは、予算統制、lDDT-PLI による評価、
そして学習データ類似度による層別である。

異なるアーキテクチャ系統を混合する選択肢については、\citet{morehead-2025-assessing} が
PoseBench において「深層学習の協調フォールディングは同種のドッキング基準手法を概ね
上回る」と報告している。合意ベースのモデルランキング自体は確立された手法であり、
\citet{kagaya-2025-structure} は CASP16 のタンパク質複合体予測で1位を得ている。

生成モデルには Boltz-2 \citep{passaro-2025-boltz} を用いた。学習カットオフが
2023-06-01 と他手法より2年新しいため、類似度は Boltz-2 自身のカットオフで計算された
列を用いる必要がある。

\section{背景：25枠という予算}

Runs N' Poses の各手法は、系ごとに5シード $\times$ 5拡散サンプル $=$ 25構造を生成する。
本研究はこの25という数を予算とみなし、その内訳のみを操作する。

評価指標は lDDT-PLI である。成功の定義は著者に合わせ、
lDDT-PLI $> 0.8$ かつリガンド RMSD $< 2$~\AA{} の連言とした。lDDT-PLI 単独では
「接触は概ね合っているが配座が違う」ポーズを成功と数えてしまうためである。

\section{手法}

\subsection{3つの配分}

総数を25に固定したまま、内訳のみを変える3アームを定義する。

\begin{itemize}
\item \textbf{baseline}：1本の MSA から25拡散サンプル。Runs N' Poses の手続きの再現。
\item \textbf{proposed}：5つの MSA 深度変種 $\times$ 各5サンプル。深度は
      100\%, 50\%, 25\%, 12.5\%, 6.25\%。変種0は完全なアライメントであり、
      \textbf{baseline の入力を部分集合として含む}。
\item \textbf{lineage}：Boltz-2 で20構造、その最高信頼度の受容体に対して DiffDock で
      5ポーズ。異なるアーキテクチャ系統への配分。
\end{itemize}

モデルの重みは変更せず推論のみに用いる。学習も微調整も行わない。

\subsection{交絡の除去}

3アームは\textbf{同一の MSA ファイル}を読む。MSA 検索は系ごとに一度だけ実行し、
結果をリポジトリに保存して全アームがそれを参照する。したがって MSA の同一性は
推定ではなく保証されている。サンプリングパラメータ（recycling 3、sampling steps 50）も
全アームで固定した。

\subsection{採点}

採点は OpenStructure の \texttt{LDDTPLIScorer} と \texttt{SCRMSDScorer} に統一した。
自前実装を用いなかったのは、\texttt{add\_mdl\_contacts=True}（OST の既定であり
plinder.eval も使用）を実装しない素朴な実装が、モデル側の作る偽の接触を数えないために
退化した予測に対して系統的に甘くなるためである。予備検証では最大0.31の乖離を確認した。
設定は plinder.eval に合わせ \texttt{substructure\_match=True}、
\texttt{lddt\_pli\_radius=6.0}、閾値 $[0.5, 1.0, 2.0, 4.0]$ とした。

\subsection{事前登録した判定基準}

データ取得前に、主指標 oracle@25 lDDT-PLI の対応差について\textbf{閾値 $+0.030$}を
凍結した。系を単位とした1,000回のブートストラップ95\%区間が0を跨いだ場合は
有意差なしと判定する。

\section{実験設定}

母集団は Runs N' Poses のうち、Boltz-2 自身の2023カットオフ類似度
（\texttt{sucos\_shape\_pocket\_qcov\_2023}）で 20--60 ビンに入る単一リガンド系とした。
20未満を除くのは全手法が一様に失敗して配分の差が現れないため、60以上を除くのは
記憶により既にほぼ解けており天井効果が生じるためである。

選択条件を満たすのは221系であった。うち1系は非標準残基（X）を含み MSA 検索が
拒否するため選択時に除外し、220系を対象とした。実行環境は理化学研究所の
GB200 クラスタ（aarch64）で、Boltz-2・DiffDock・MSA 検索はいずれも
NVIDIA NIM のホスト型エンドポイントを呼ぶ。

\section{結果}

\subsection{ランキングが取り逃している量：ベースラインでの再現}

baseline において、25構造中の最良（oracle@25）は 0.6213、そのアーム自身の信頼度で
選んだ1本（top-ranked）は 0.5560 であった。\textbf{差は 0.065} である
（以下、本稿の数値はすべて3アーム共通の205系上のものである）。

これは Runs N' Poses が AlphaFold3 について報告した $0.830 - 0.774 = 0.056$ と
同じオーダーである。\textbf{ただし両者を直接比較することはできない。}著者の値は全類似度帯
2,292系における AlphaFold3 の値であり、本研究の値は 20--60 帯205系における Boltz-2 の
値である。母集団も対象モデルも異なるため、言えるのは「別のモデルと別の部分集合でも
同程度の取り逃しが観測される」ところまでである。いずれにせよ、候補集合の中に当たりは
存在するのに信頼度スコアがそれを選べていない、という構図は共通している。

\subsection{主指標：閾値に届かなかった}

3アーム共通の205系での対応比較を表~\ref{tab:main} に示す。

\begin{table}[htbp]
\centering
\caption{3アームの対応比較（共通205系）。対応差はいずれも baseline を基準とし、
括弧内は系単位1,000回ブートストラップの95\%区間。}
\label{tab:main}
\begin{tabular}{lccc}
\toprule
指標 & baseline & proposed & lineage \\
\midrule
oracle@25 lDDT-PLI & 0.6213 & 0.6335 & \textbf{0.6511} \\
\quad 対応差 & --- & $+0.0121$ & $+0.0298$ \\
\quad 95\%区間 & --- & $[+0.0014, +0.0226]$ & $[+0.0187, +0.0425]$ \\
\quad 閾値 $+0.030$ & --- & 未達 & 未達（$0.0002$ 差） \\
\midrule
top-ranked lDDT-PLI & 0.5560 & 0.5391 & 0.5479 \\
\quad 対応差 & --- & $-0.0168$ $[-0.0333, -0.0017]$ & $-0.0081$ $[-0.0210, +0.0049]$ \\
\midrule
リガンド RMSD (\AA) & 4.201 & 3.964 & \textbf{3.332} \\
\quad 対応差 & --- & $-0.302$ $[-0.514, -0.118]$ & $-0.870$ $[-1.352, -0.459]$ \\
\midrule
success rate$^{\dagger}$ & 0.3659 & 0.3854 & \textbf{0.4098} \\
モード数（平均） & 3.76 & 5.21 & \textbf{6.99} \\
\bottomrule
\end{tabular}

\vspace{1mm}
{\footnotesize $\dagger$ success rate は系ごとの二値であり、対応差が $\{-1,0,+1\}$ しか
取らない。他の連続指標と同じ形の区間を併記していないのはそのためである。}
\end{table}

\paragraph{proposed。}対応差は $+0.0121$ で95\%区間は0を跨がず、効果の存在は
否定されない。\textbf{しかし事前に登録した閾値 $+0.030$ には届かなかった。}区間の上限
$+0.0226$ すら閾値を下回る。改善した系は121、悪化した系は79であった。

\paragraph{lineage。}対応差は $+0.0298$、区間 $[+0.0187, +0.0425]$。区間は0を大きく
上回り、上限は閾値を超える。\textbf{しかし凍結した判定規則は平均が $+0.030$ 以上で
あることを求めており、平均 $+0.0298$ はこれを $0.0002$ 下回る。}したがって規則に従えば
\textbf{未達}である。四捨五入すれば閾値に一致する値であり、判定境界のほぼ真上に
乗っているが、事前に凍結した規則を事後に緩めることはしない。この近さ自体が、
閾値を1点で切ることの限界を示す所見でもある。

\subsection{効果は難しい系ほど大きい、わけではない}

類似度ビン別の対応差を図~\ref{fig:bin} と表~\ref{tab:bin} に示す。

\begin{table}[htbp]
\centering
\caption{類似度ビン別の oracle@25 対応差（対 baseline、共通205系）。
太字は区間が0を跨がないもの。}
\label{tab:bin}
\begin{tabular}{lccc}
\toprule
ビン & $n$ & proposed & lineage \\
\midrule
20--30 & 27 & $-0.0170$ $[-0.0487, +0.0054]$ & \textbf{$+0.0270$} $[+0.0079, +0.0511]$ \\
30--40 & 40 & $+0.0367$ $[-0.0027, +0.0786]$ & \textbf{$+0.0412$} $[+0.0088, +0.0817]$ \\
40--50 & 50 & \textbf{$+0.0218$} $[+0.0088, +0.0419]$ & \textbf{$+0.0428$} $[+0.0161, +0.0783]$ \\
50--60 & 88 & $+0.0086$ $[-0.0020, +0.0208]$ & \textbf{$+0.0181$} $[+0.0070, +0.0313]$ \\
\bottomrule
\end{tabular}
\end{table}

\begin{figure}[htbp]
\centering
\includegraphics[width=0.72\textwidth]{images/chart/oracle_by_bin.pdf}
\caption{類似度ビン別の oracle@25 lDDT-PLI。誤差棒は系単位ブートストラップ95\%区間。}
\label{fig:bin}
\end{figure}

proposed について、我々は事前に「効果が存在するのは学習データから遠い難しい系」と
予想していた。\textbf{データはこれを支持しない。}最も難しい 20--30 帯で対応差は負に転じ、
区間が0を跨がないのは中間の 40--50 帯だけである。

\textbf{lineage は対照的に、4つのビンすべてで区間が0を跨がない。}効果は難易度に
依存せず一貫している。

proposed について素直な解釈は、浅い MSA が両刃だということである。学習データに近い系では
共進化信号に冗長性があるため深度を落としても失うものが少なく、多様化の効果だけが残る。
一方 20--30 帯では信号そのものが乏しく、深度を落とすと\textbf{頼るべき情報まで削って
しまう}。ただし当該ビンの区間は0を跨ぐため、悪化を断定はできない。断定できるのは
「難しい系ほど効く」という当初の期待が支持されないことまでである。

この非対称性は、2つの配分が触っているものの違いとして理解できる。MSA 摂動は
\textbf{同じモデルの入力}を揺らすため、その効果は入力に含まれる情報量に依存する。
別系統への配分は\textbf{異なる分布から標本を引く}操作であり、元の分布が何であれ
成立する。

\subsection{摂動は作動していた：機構確認}

否定的な主結果は、摂動が弱すぎたためではない。図~\ref{fig:modes} に、25構造を
対称性補正済みリガンド RMSD 2~\AA{}・完全連結でクラスタリングしたときのモード数を示す。

\begin{figure}[htbp]
\centering
\includegraphics[width=0.72\textwidth]{images/chart/pose_modes.pdf}
\caption{25構造中の区別できるポーズモード数の分布。平均 3.77 → 5.21。}
\label{fig:modes}
\end{figure}

平均モード数は baseline 3.76、proposed 5.21、lineage 6.99 である。モードが1本しか
ない系は baseline 69、proposed 55、\textbf{lineage 6} と大きく減る。
すなわち\textbf{別系統の混合は MSA 摂動より強く多様性を生んでいる}。

この増加が採点失敗による水増しでないことは別途確認した。スコアできないポーズ対は
クラスタに併合されないため必ずモード数を押し上げるが、その割合は baseline 0.0246 に対し
proposed 0.0308 と差が0.6ポイントしかなく、この増加を説明できない。

したがって、事前に用意した2つの説明のうち\textbf{「摂動が弱すぎた」は棄却される}。
残るのは\textbf{「多様性は増えたが、増えた方向が正解の方向ではなかった」}である。

\subsection{予期しなかった所見：候補を広げると選択が悪化する}

表~\ref{tab:main} と図~\ref{fig:gap} が示すとおり、proposed では oracle が上がった
一方で\textbf{top-ranked は有意に下がった}（対応差 $-0.0168$、区間
$[-0.0333, -0.0017]$）。lineage でも同じ向きだが、こちらは区間が0を跨ぐため
有意な悪化とは言えない（$-0.0081$、$[-0.0210, +0.0049]$）。

oracle と top-ranked の差、すなわちランキングが取り逃す量は、
\textbf{0.065（baseline）$\to$ 0.094（proposed）$\to$ 0.103（lineage）と
候補の多様性に伴って単調に増加する}。モード数の平均が
3.76 $\to$ 5.21 $\to$ 6.99 と増える順序と一致している。

\begin{figure}[htbp]
\centering
\includegraphics[width=0.62\textwidth]{images/chart/oracle_vs_topranked.pdf}
\caption{oracle@25 と top-ranked の対比。候補集合を広げると両者の差が開く。}
\label{fig:gap}
\end{figure}

リガンド RMSD は有意に改善している（$-0.302$~\AA、区間が0を跨がない）にもかかわらず、
信頼度で選んだ1本は有意に悪化している。\textbf{幾何学的により近いポーズを生成できる
ようになっても、それを選び出せていない。}

これは Runs N' Poses の中核の観察——信頼度は系の難易度を測るが系内のどれが正しいかは
測らない——を一段強める。予算配分という枠組みで言い換えれば、\textbf{多様性を買うことは
選択の難しさを買うことでもある}。

\subsection{図：系ごとの対応差}

\begin{figure}[htbp]
\centering
\includegraphics[width=0.72\textwidth]{images/chart/paired_differences.pdf}
\caption{系ごとの対応差の分布。実線が0、破線が事前登録した閾値 $+0.030$。
分布の質量の大半が閾値の左側にある。}
\label{fig:paired}
\end{figure}

\subsection{外れた事前予測：別系統への配分}
\label{sec:lineage}

我々はこのアームについて事前に\textbf{負の結果を期待値}として登録していた。
\citet{morehead-2025-assessing} が深層学習の協調フォールディングは同種のドッキング
基準手法を概ね上回ると報告していること、および拡散ドッキングが予測受容体に対して
劣化することが知られているためである。\textbf{この予測は外れた。}

lineage は3アーム中で最良の oracle@25（0.6511）を示し、リガンド RMSD は
$0.87$~\AA{} 改善した（区間 $[-1.352, -0.459]$）。成功率も 0.3659 から 0.4098 へ
上がっている。

\paragraph{なぜ効いたのか。}ドッキングは25枠のうち5枠、すなわち予算の20\%を占める。
そして\textbf{25本中の最良ポーズがドッキング由来であった系は 41/205 = 20.0\%}であった。
両者はほぼ完全に比例する。すなわち\textbf{ドッキングのポーズは1本あたりの質において
協調フォールディングと互角であり、優れているわけではない。}

にもかかわらず oracle が上がるのは、最大値の統計としての性質による。同一の分布から
25本引くより、相関の低い2つの分布から20本と5本を引く方が、最大値の期待値は高くなる。
\textbf{効いているのはポーズの質ではなく、分布の非相関性である。}これは
「独立に学習された系統は異なる誤り方をする」という当初の想定と整合するが、
我々はその効果が既知の性能差に打ち消されると予想していた。打ち消されなかった。

\section{制約}

\textbf{予算が等しいのは構造数であって計算時間ではない。}proposed は1系あたり平均
5.70回の API 呼び出しを要し、baseline の1.47回より多い（baseline が1回を超えるのは、
大きな複合体で25本一括のリクエストが失敗した際に分割して同数を生成したためである）。
lineage は協調フォールディングに加えてドッキングの呼び出しを要する。科学的に意味のある
単位は構造数だと考えるが、実運用のコストはこの差を含む。

\textbf{母集団は220系のうち205系である。}3アームすべてで採点できた系のみを用いた。
脱落の主因は生成基盤の障害窓であり、そのうち生成器が25本を作れなかった系の失敗モードは
系のサイズと相関するため、\textbf{脱落は完全にランダムではない}。ただし対応比較は
3アーム共通の集合上で行うため、アーム間の非対称な脱落は生じない。

\textbf{摂動の種類は1つしか試していない。}本研究が検証したのは深度サブサンプリング
のみであり、カラムマスクやクラスタリングは未検証である。したがって否定的結論は
\textbf{深度サブサンプリングについてのもの}であって、MSA 摂動一般についてではない。

\textbf{単一のモデル対での結果である。}協調フォールディングは Boltz-2、ドッキングは
DiffDock のみを用いており、他の組合せで同じ結論が成り立つかは検証していない。

\textbf{lineage の受容体は固定していない。}ドッキング先は協調フォールディングが出した
最高信頼度の受容体であり、その品質が結果に影響する。結晶構造の受容体を用いた上限参照は
本研究では取っていない。

\section{結論}

25構造という固定予算の配分について、3つの戦略を対比較した。

\textbf{MSA 深度のサブサンプリングは、事前に登録した閾値に届かなかった。}摂動自体は
確実に作動し、ポーズの多様性を約4割増やしたにもかかわらず、oracle@25 の改善は
$+0.0121$ にとどまった。効果は中間的な難易度の系に限られ、最も難しい系ではむしろ
負に転じた。したがって、\textbf{上流のタンパク質側を揺らす操作は、外している当の
リガンド配座には届かない}。これは Runs N' Poses の著者が示した「推論時に複数シードは
有用でない」と並ぶ、サンプリング側の2つ目の否定的データ点である。

\textbf{別のアーキテクチャ系統への配分は、閾値を $0.0002$ 下回った。}平均 $+0.0298$、
区間 $[+0.0187, +0.0425]$。判定規則に従えば未達だが、4つの類似度帯すべてで区間が
0を跨がず、リガンド RMSD は $0.87$~\AA{} 改善した。\textbf{我々はこのアームについて
負の結果を事前登録しており、予測は外れた。}

機構は明快である。ドッキングは予算の20\%を占め、最良ポーズを供給した割合も 20.0\%で
比例していた。すなわち1本あたりの質は互角であり、\textbf{効いているのは分布の
非相関性}である。同一分布から25本引くより、相関の低い2つの分布から引く方が最大値が
高くなる。

以上を合わせると、25枠の配り方についての答えは\textbf{「同じモデルの入力を揺らす」
のではなく「別の分布から引く」}方向にある。ただしどちらも事前に設定した意味のある
大きさには届いておらず、\textbf{配分の変更だけで解ける問題ではない}というのが
本研究の最も正直な要約である。

最後に、3アームを通じて一貫した所見として、候補集合が広いほど信頼度による選択が
悪化した。サンプリングを改善してもランキングがそれを拾えなければ実務上の利得に
ならない。予算をサンプリングに投じる前に、選択器そのものを改善する方が利得が大きい
可能性がある。

本研究が主張するのは\textbf{固定推論予算下での配分戦略の比較}であって、新しい協調
フォールディングモデルではない。

\paragraph{データと再現性。}全アームが参照した MSA、系ごとの全ポーズのスコア、
および解析コードは実験リポジトリに保存されている。AlphaFold3 の出力を含む
Runs N' Poses のデータには DeepMind の Output Terms of Use が付随する。

\bibliography{references}

\end{document}
